\section{Replication、Cluster Formation 与 Hierarchical Sparsity}

GLOM 会在一个 object 覆盖的许多位置复制同一高层 embedding，看起来非常浪费。Hinton 的反驳是：在 binding 尚未确定时，每个位置保留独立 hypothesis 才能逐步决定“哪些位置应该相同”；复制不是最终压缩最优，而是 inference flexibility 与 locality 的代价。

\lecturefigure{slide-30-object-embedding-replication.jpg}{每个 object patch 复制 object-level embedding，看似浪费却保留局部 hypothesis。}{Stanford Online 官方视频 00:48:52--00:49:47。}

\teachervoice{讲者用 biology 类比：细胞复制相同 DNA，器官不同部位也可有相近 protein-expression vector。这个类比的教学作用是说明 locality 可能值得冗余；它不是说 GLOM embedding 等同于基因或蛋白表达。}

\subsection{Forming clusters，不是对固定点做 clustering}

普通 clustering 接收固定 data points，再发现分组；GLOM 中 data points 自己会随 bottom-up、top-down、lateral 与 temporal interaction 改变。系统不是对既有 embedding 聚类，而是在 recurrent dynamics 中共同形成 embedding 与 cluster boundary。

\begin{importantbox}{“Hedge your bets”的计算含义}
在早期 step，位置 $x$ 可以保留与邻居不同的 parent hypothesis；随着证据积累，兼容 positions 才被吸入同一 island。若一开始就把整块 object 压成单一 shared variable，错误 binding 很难局部修正。复制换来的是 delayed commitment。
\end{importantbox}

\begin{warningbox}{动态形成不等于自动获得正确 discrete cluster}
连续向量可能形成模糊过渡、多个局部 attractor 或 global consensus。若需要离散 object slot，还要定义 readout、threshold、connectedness、temporal tracking 与 birth/death rule。本讲并未完成这套 contract。
\end{warningbox}

\subsection{高层更远、更稀疏}

低层 part 小、边界细，需要短程 dense interaction；高层 object/scene island 大、内部冗余高，可以扩大 receptive radius，却只采样少数位置。目标是让 hierarchy level 上升时，communication range 增长而 edge count 不爆炸。

\lecturefigure{slide-31-sparse-high-level-replication.jpg}{高层 islands 更大，可使用更长程、更稀疏的 attention connection。}{Stanford Online 官方视频 00:49:48--00:50:07。}

\begin{knowledgebox}{读图：range 增大不要求 edge count 同比例增长}
高层 object island 覆盖许多 patches，内部向量预期高度冗余，因此一个位置不必连接 island 内所有位置；只要 sparse sample 命中代表点，就可能获得相同 whole hypothesis。图表达的是 range--density tradeoff：interaction radius 随 level 增长，连接密度下降。它没有给出 neighbor-search cost、miss rate 或实际吞吐量，所以只能作为 systems hypothesis。
\end{knowledgebox}

\begin{table}[H]
\centering
\small
\caption{GLOM hierarchy 中 locality、range 与 sparsity 的设计意图。}
\begin{tabularx}{\textwidth}{L{0.17\textwidth}L{0.24\textwidth}L{0.24\textwidth}>{\raggedright\arraybackslash}X}
\toprule
层级 & 典型实体 & Interaction pattern & 主要风险 \\
\midrule
低层 & edge、texture、small part & 短程、较 dense，保护细边界 & 只看局部会错过 whole context \\
中层 & major part、object component & 中程、按 similarity gated & 容易把相似但不同实例粘连 \\
高层 & object、person、scene & 长程、较 sparse，利用大 island 冗余 & sparse sample 可能漏掉小对象或多实例冲突 \\
\bottomrule
\end{tabularx}
\end{table}

\begin{warningbox}{“每层工作量相同”是目标，不是测量结果}
要验证该主张，需要明确每层位置数、neighbor count、embedding width、迭代步数、routing/ANN search、memory traffic 与 communication pattern。单说“更 sparse”不足以证明 FLOPs、latency 或 energy 恒定。
\end{warningbox}

\subsection{本章小结}

Embedding replication 为 gradual binding 与 local correction 保留自由度；hierarchical sparsity 试图把这种冗余转化为可扩展 interaction。两者共同体现 GLOM 的核心取舍：宁可在 state memory 上重复，也不提前把 uncertain structure 压成不可逆的单一 node。

